\section{Introduction}
Archaeological conservation requires decisions across spatial scales: monument fabric, its immediate surroundings and the wider landscape. Satellite observations, terrain models and weather products capture different parts of this system. Translating them into an inspection programme requires more than overlaying indicators: the resulting priorities must be traceable to their environmental drivers and stable enough for a defined management task.

Remote sensing and GIS have documented development pressure and supported heritage multi-hazard assessment \citep{agapiou2015paphos,agapiou2016paphos,guerriero2022derwent}. At Taxila, Khan et al. integrated ground-based and satellite evidence for hazard assessment \citep{khan2022taxila}; Butt et al. subsequently examined land-cover transitions and urban encroachment over 1990--2024 \citep{butt2025taxila}. These studies establish the local conservation problem. Recent work by Megarry et al. extends the relevance of surrounding land cover to climate-vulnerability assessment across cultural World Heritage properties in the Indian Subcontinent \citep{megarry2026landcover}. The next decision-oriented question is whether an inspection set persists when spatial support, indicator composition and decision weights change. Mapping pressure and evaluating the stability of a management recommendation are complementary contributions.

GIS multicriteria methods formalise indicator aggregation, while sensitivity analysis reveals dependence on weights and spatial representation \citep{malczewski2006gis,feizizadeh2014uncertainty}. Rank-acceptability methods further describe which alternatives remain preferred under uncertain criteria and preferences \citep{lahdelma2001smaa}. For archaeological landscapes, these ideas motivate an experimental treatment of priority construction: correlated proxies can duplicate information, neighbouring buffers can share observations, and uncertain coordinates can alter the sampled environment. A single final ranking leaves these dependencies unresolved.

Climate interpretation introduces temporal-support constraints. Exposure depends on materials, maintenance, drainage and microenvironment alongside meteorological forcing \citep{sesana2021climate,sesana2020vulnerability,crowley2022risk}. Acquisition-matched weather and cross-product comparison can test whether a proposed explanation is consistent with the observations. Separately, predictive evaluation must distinguish geographic land-cover transfer from conservation-condition assessment. Spatially nearby samples can share predictors and labels \citep{hijmans2012cv,roberts2017cv,meyer2021aoa}; reproducing a global land-cover product does not establish monument deterioration.

We develop CHIP as an experimentally evaluated inspection framework with three contributions. First, a hierarchical landscape--terrain rule exposes indicator influence through baselines and ablations, and exact weight crossings identify where inspection leadership changes. Second, multiscale, shared-block, positional and nested spatial--decision experiments distinguish persistent priorities from representation-dependent ones. Third, acquisition-matched climate comparisons and a seven-family buffered land-cover benchmark evaluate the environmental and predictive evidence supporting interpretation. The resulting inspection portfolio retains both a persistent component set and the mechanisms that require field verification.

\section{Research aim}
The study identifies which Taxila components merit early field inspection and which parts of that recommendation depend on analytical choices. It tests adverse spectral convergence on common spatial support, changes in climate interpretation across acquisition windows and products, rank stability across neighbourhoods and uncertainty designs, and land-cover proxy agreement under buffered spatial evaluation. Mathematical definitions connect these experiments to reproducible component-level decisions.
